%% kapa: exam=1 round=33 year=2022 period=1 form=A subject=civil_law subject_ko=민법 seq=35 number=35
%% source: Q-Net 감정평가사 기출문제 (cst00309) artlSeq=5209997 file=raw/1008/5209997/2022년도 제33회 감정평가사 1차시험 1교시(민법,경제학,부동산학).pdf
%% conversion: kordoc 4.24.51 -> tools/gen1.py (auto); review=proofread-sonnet; flags=none
\begin{question}{35}
민사유치권에 관한 설명으로 옳은 것은? (다툼이 있으면 판례에 따름)
\begin{choices}
  \item 유치권 배제 특약이 있더라도 다른 법정요건이 모두 충족되면 유치권이 성립한다.
  \item 채무자는 상당한 담보를 제공하고 유치권의 소멸을 청구할 수 있다.
  \item 원칙적으로 유치권은 채권자 자신 소유 물건에 대해서도 성립한다.
  \item 채권자가 채무자를 직접점유자로 하여 간접점유하는 경우, 채권자의 점유는 유치권의 요건으로서의 점유에 해당한다.
  \item 채권자의 점유가 불법행위로 인한 경우에도 유치권이 성립한다.
\end{choices}
\end{question}
\answer{2}
